\section{Discussion}\label{sec:discussion}

\subsection{What the results establish}

Calibrating the alarm rather than the interval did what Theorem~\ref{thm:main} says it does. On seven datasets, with between 32 and 83 calibration units, the failure-before-replacement rate followed its nominal level, while a per-cycle lower bound at the same nominal level, on the same signal and trigger, failed on up to 70\,\% of units; the contrast held with four learners and the guarantee with 27 alarm settings. The difference is not a matter of tuning: per-cycle coverage averages over cycles chosen independently of the data, whereas an alarm is read at the first cycle at which the signal crosses a level, the cycle at which an optimistic error is most likely. Restricting the per-cycle residuals to the near-failure region repaired most of the gap for GBM on six datasets, but how close it came depended on that choice of rows and on the learner: on N-CMAPSS it exceeded its level by half, and with ridge regression it failed on 17--33\,\% of units. Calibrating the critical level removes the choice, because the quantity calibrated is the event the operator cares about. The practical message for the growing use of conformal RUL intervals in maintenance~\cite{javanmardi2023,robinson2026,wang2025robustuq} is that interval coverage should not be read as alarm reliability.

\subsection{A deployment recipe}

Table~\ref{tab:recipe} turns the results into a procedure. Its steps are the ones that, in this study, changed whether the stated reliability was achieved; none requires a model of the degradation or of the RUL error.

\begin{table}[t]\footnotesize\setlength{\tabcolsep}{3pt}
\caption{Deployment recipe for a notice-guaranteed end-of-life alarm.}\label{tab:recipe}
\begin{tabular}{@{}p{0.04\linewidth}p{0.93\linewidth}@{}}
\toprule
& Step and reason \\
\midrule
1 & State the lead time $\ell$ and the tolerated probability $\alpha$ of failure before the planned replacement; check $n\ge1/\alpha-1$ calibration units are available (Proposition~\ref{prop:cost}). \\
2 & Fix the signal and trigger (smoothing, debounce, cap) before looking at the calibration units, or select them on units not used for calibration. \\
3 & Collect calibration units from the population in which the alarm will run. Use replaced or still-running units as censored units, their critical level computed as if they had failed when last seen alive (Theorem~\ref{thm:cens}). \\
4 & Compute critical levels (Lemma~\ref{lem:crit}) and set $\hat H$ to their conformal quantile; with a cost ratio, choose among admissible order statistics (Corollary~\ref{cor:da}). \\
5 & When the population changes (new batch, campaign or operating regime), recalibrate on units of the new population, starting from a conservative interim level; do not carry a calibrated level across populations. \\
6 & Report the failure rate with its interval, the notice given, and the number of calibration units the claim rests on; do not report interval coverage or a window score as alarm reliability. \\
\bottomrule
\end{tabular}
\end{table}


\subsection{What the guarantee costs}

The guarantee is a certificate; whether it also saves money depends on how many calibration units it rests on. Proposition~\ref{prop:cost}(ii) gives the mechanism: a level chosen among $n$ calibration critical levels fails with probability at least $1/(n+1)$, and with expensive failures that floor dominates the cost. With the split design and 32 calibration units the certified alarm was clearly dearer than the best uncertified rules; cross-fitting, which calibrates on all training units and trains the signal on four fifths of them, removed most of the premium (Table~\ref{tab:cost}). With equal information, cross-fitted calibration made the certificate almost free: the certified alarm was never significantly dearer than the best uncertified rules built on the same models, and cheaper than the one-standard-error rule on two datasets (Section~\ref{sec:res-cost}).
 Rules that choose levels beyond the most demanding calibration unit can be cheaper on a given test set, and the earlier study showed how such rules can fail when plain empirical cost minimisation finds a setting at the edge of the feasible region~\cite{prior}; what they cannot provide is a bound for the next unit. Two sizing rules follow. To certify a failure probability $\alpha$ one needs at least $1/\alpha-1$ run-to-failure (or late-censored) calibration units from the population in question; to keep the certified excess cost $(\rho-1)/(n+1)$ below a fraction $\varepsilon$ of a planned replacement one needs $n+1\ge(\rho-1)/\varepsilon$. Where such data exist, as for fleets of identical components or cells from one production line, the certificate is cheap; where they do not, a failure rate of an uncertified alarm measured on a few dozen units is weak evidence of its reliability.


\subsection{Transfer and censoring in the field}

The battery results show the limit of every exchangeability-based guarantee: calibrated on two batches, the alarm failed on most cells of a third, shorter-lived batch, and similar failures occurred on three of nine N-CMAPSS subsets. This is not a defect of the conformal step, which did what it promised for the population it was calibrated on, but it is a warning against reading a validated failure rate as a property of the alarm rather than of the alarm and a population. Three responses were examined. Conservative uncertified rules carried enough margin to survive this shift, at the price of notice that is wasted when the population does not change. Recalibration on ten pilot cells of the new batch restored the per-unit guarantee, but those cells must be run close to failure, because a cell replaced under a conservative level can confirm that level but never relax it. Online calibration (Theorem~\ref{thm:online}) needs neither exchangeability nor pilot failures: it uses only whether each new unit failed before its replacement, which is observed in any fleet, and it brought the failure rate of the shifted batches back to about 10\,\% after the first half of each sequence. Its guarantee is of a different kind, a bound on the realised long-run rate rather than a probability for the next unit, and short sequences, early failures in a strongly shifted population and delayed feedback limit it. Weighted or non-exchangeable conformal methods~\cite{tibshirani2019,barber2023} could use covariates of the new batch, such as its charging protocol, to reweight calibration units; this needs a model of the shift and was not attempted.

\subsection{Limitations}

\begin{itemize}[leftmargin=1.2em,itemsep=2pt]
\item \textbf{Data.} Five of the seven datasets come from one simulator family and the battery cells are laboratory data. Censoring was imposed on complete run-to-failure records by three mechanisms; no field fleet with genuinely censored lives was analysed, and none was available to us, nor was a second laboratory dataset of the same cell type that would test transfer across laboratories. Bearings were excluded because 17 units cannot support the guarantee.
\item \textbf{Exchangeability.} The per-unit guarantees require exchangeable calibration and new units; they failed across battery batches and N-CMAPSS subsets and may fail under slower drifts within one population. The online update replaces them by a bound on the realised long-run failure rate that holds for any sequence but says nothing about an individual unit, is weak for short sequences and assumes that outcomes are fed back promptly. Theorem~\ref{thm:main} is marginal over calibration sets; the training-conditional form (Corollary~\ref{cor:pac}) needs independent units, and making it hold with high confidence costs notice.
\item \textbf{Scope.} The theory covers threshold triggers with debounce on a causal signal and a fixed lead time. The settings of CNA (cap, smoothing, debounce) were fixed a priori, not optimised; 27 alternative settings gave pooled failure rates of 8.7--11.1\,\% and almost the same notice. The cross-fitted guarantee (Theorem~\ref{thm:cvplus}) is weaker than the split guarantee.
\item \textbf{Cost.} The renewal--reward model assumes independent units, replacement as good as new and illustrative cost ratios. Section~\ref{sec:res-m} uses the decision setting of Kamariotis et al.\ but our own implementation of their policy and of the probability model behind it.
\item \textbf{Pre-specification.} The protocol and its four amendments were hashed before their runs but not registered, and each amendment was written after earlier results were known. The equivalence and cost hypotheses of the main protocol and the cost hypothesis of the first amendment (H7b) failed and are reported as such; the prominence of Proposition~\ref{prop:cost}(ii) as the explanation of the cost results was decided after those results were seen, and the check of Corollary~\ref{cor:pac} was added after all runs.
\end{itemize}

\section{Conclusion}\label{sec:conclusion}

% stub sec9_conclusion


\section*{Declaration of generative AI and AI-assisted technologies in the manuscript preparation process}

\noindent During the preparation of this work the authors used Claude (Anthropic), an AI assistant, to search the literature, write and run the analysis code, draft the theory and proofs, and draft and revise the manuscript text. The protocols, all experiments and all reported numbers were produced by the code in the accompanying repository and checked against its output files. After using this tool the authors reviewed and edited the content as needed and take full responsibility for the content of the published article.

\section*{Data and code availability}

\noindent PHM08, C-MAPSS and N-CMAPSS are available from the NASA Prognostics Center of Excellence data repository, and the battery data of Severson et al.~\cite{severson2019} from \url{https://data.matr.io}. Code, hashed protocols, derived data summaries, per-unit results and the scripts that produce every table and figure are in a repository whose link is withheld for anonymised review.
